

\documentclass[11pt]{article}
\usepackage[margin=1in]{geometry}
\usepackage{amsmath,amssymb,amsthm}
\usepackage{graphicx}
\usepackage{booktabs}
\usepackage{enumitem}
\usepackage{microtype}
\usepackage[hypertexnames=false]{hyperref}
\usepackage{algorithm}
\usepackage{algpseudocode}

\title{Advantage Contracts for Anonymous DHT Evaluations:\\
Stealth Audits, Leakage Discovery at Scale, and Practical Padding Filters}
\author{}
\date{March 19, 2026 (archive v393)}

\newtheorem{theorem}{Theorem}
\newtheorem{lemma}{Lemma}
\newtheorem{proposition}{Proposition}

% Shared notation/style across the Anonymity / Anonymous DHT paper series.

\newcommand{\TV}{\mathrm{TV}}
\newcommand{\E}{\mathbb{E}}
\newcommand{\Pp}{\mathbb{P}}

% Reduce overfull boxes without forcing full \sloppy.
\setlength\emergencystretch{2em}

\begin{document}
\maketitle

\begin{abstract}
Anonymous DHTs and overlay lookups are typically evaluated with active probes: canary lookups, synthetic workloads, and measurement traffic.
A practical adversary can \emph{detect} these audits and behave differently under evaluation than in deployment (``sandbagging''),
biasing published anonymity and reliability claims.
We propose an \emph{advantage contract} that makes evaluation bias explicit and controllable:
bound the audit--live distinguishability (TV advantage) of the observable traces used by the evaluator.
Under this contract, any bounded evaluation metric (risk, leakage proxy, success rate) cannot shift by more than the advantage budget.
We then make the contract operational with (i) anytime-valid monitoring of detectability via e-processes, (ii) scalable slice discovery with e-values and e-BH,
and (iii) practical stealth mechanisms (commit--reveal probe selection with public randomness) plus padding filters that minimize overhead subject to an advantage constraint.
We include a lightweight Kademlia-flavored trace case study showing how detectable audit infrastructure induces audit-live bias and how stealth reduces both.
\end{abstract}

\paragraph{Compilation note.}
This paper is brief-by-design. For extended deployed-evidence pointers and additional bibliography, see the related work map.\cite{series:relatedwork}

\paragraph{Scope.}
We study \emph{evaluation awareness}: the system under test may condition on whether traffic is an audit, inducing an audit--live distribution shift.
We package this shift as a published \emph{advantage contract} that upper-bounds audit detectability on the released trace schema.
This paper owns only that audit--live distinguishability layer and the declared stealth-bias packet attached to it; it does not own runtime conformance to a pinned plan (Operational~B), evaluator notarization / retry logging (Eval~1), or cross-release lineage / publication packaging (Release~A).

\paragraph{Imports.}
We import the trace-schema / exposure-tier vocabulary and the TV/padding compiler toolchain from the companion series notes.\cite{series:schedcompiler,series:knapsacktransport}
For stopping-time leakage and deadline adaptivity, we import the stop-time addendum.\cite{series:stoptimeaddendum}
For certificate mechanics and retry discipline, we import the notarized certificates note.\cite{series:evalnotary}
Exposure-normal-form identity and interface/coarsening semantics are imported from the trace-interface note.\cite{series:traceifaces}

\paragraph{Exports.}
We export (i) an advantage contract (TV distinguishability) linking audit detectability to bounded bias in any reported metric, (ii) anytime-valid monitoring of detectability via e-processes and scalable slice discovery (e-values, e-BH), and (iii) practical stealth mechanisms (commit--reveal probe selection with public randomness) and padding filters that reduce detectability subject to an advantage budget.
The stable public object is a compact \emph{stealth-bias packet}: released trace schema, audit/live sampling declaration, declared audit-advantage bound (or confidence interval), slice-discovery correction rule, and any declared stealth mechanism / padding filter.
This paper deliberately stops before evaluator notarization and retry-prefix logging (Eval~1), plan-pinning / PTL / epoch-conformance receipts (Operational~B), and receipt lineage / release-ledger comparison (Release~A).

\paragraph{Earliest sufficient stop.}
If a deployment can publish the released trace schema, the audit/live sampling rule, the declared audit-advantage bound (or interval), and the correction rule for any slice-discovery search, then the audit-awareness question answered by this paper is already fixed.
Everything after that---notarizing the attempt prefix, proving runtime conformance to a pinned plan, and packaging a lineaged release receipt---belongs to Eval~1, Operational~B, and Release~A respectively.

\paragraph{Later terse-paper citation posture.}
For later terse papers under space pressure, default to this paper for the public audit-awareness / stealth-bias claim.
There is no hidden support addendum that later papers must cite to make that claim intelligible: the exact released trace schema, audit/live sampling declaration, declared advantage budget, slice-discovery correction rule, and any declared stealth mechanism / padding filter already live here.
If the live question is whether a published trace-level metric transfers from audits to deployment up to the declared detectability budget, stop at this paper.
If the question becomes retry-prefix logging / detector-change accounting, runtime plan conformance, or public release-lineage comparison, leave this paper and reopen Eval~1, Operational~B, or Release~A respectively.
The maintained worked example and paired cutover rule are the general model for that stop / reopen discipline.\cite{series:workedexample,series:pairedterminalcutoverrules}

\paragraph{Stable handoff to neighboring layers.}
This note intentionally stops at the earliest stage that answers the audit-awareness question.
Its owned spine is:
\emph{released trace schema / audit-live sampling rule / declared advantage budget / slice-discovery correction rule}.
That spine is enough to state how much a published metric can move between audit and live traffic, but it is not enough to certify retry discipline, runtime-plan conformance, or public release lineage.
The surrounding stop map is therefore short and explicit:
\begin{itemize}[leftmargin=1.2em,itemsep=0.2em,topsep=0.35em]
\item \textbf{Audit-awareness transfer claim:} this paper's stealth-bias packet.
\item \textbf{Retries / aborts / detector changes logged under a fixed evaluator workflow:} Eval~1's manifest, gap-free attempt prefix, and verifier report.
\item \textbf{Runtime traces still matched the pinned plan:} Operational~B's plan certificate, PTL inclusion proof, and epoch conformance receipt.
\item \textbf{Comparison against earlier public releases:} Release~A's lineaged release receipt / ledger entry.
\end{itemize}
\paragraph{A minimal public advantage-contract card.}
When a later short paper only needs to answer \emph{which audit/live transfer packet was declared for this evaluation family?}, it should stop at Table~\ref{tab:adv_manifest} plus the tiny transfer drill below.

\paragraph{Artifact policy.}
This archive is source-only; supporting artifacts (datasets, plots, scripts, and tooling) are available on request.\cite{series:artifactpolicy}


\section{Series context and scope}
This paper connects to two prior themes in the series:
(1) \emph{trace-based} anonymity objectives and TV leakage bounds, and
(2) padding/scheduling compilers and transport-style certificates \cite{series:schedcompiler,series:knapsacktransport}.
We do not re-prove these foundations.
For stopping-time leakage and deadline adaptivity on the privacy side, see the series addendum \cite{series:stoptimeaddendum}.
Our focus is \emph{evaluation awareness}: the system under test may condition on whether traffic is an audit.
This is analogous to overfitting a test/holdout set under adaptive reuse: if a system can infer which queries are audits, it can ``fit'' to the evaluator.
The statistics literature studies such adaptivity and mechanisms for safe holdout reuse \cite{dwork2015reusableholdout,dwork2014adaptive}.

\paragraph{Series vocabulary (shared).}
We use a small, consistent vocabulary across the series to keep papers non-redundant.
A \emph{plan id} names an evaluation definition (trace schema, attempt unit, within-attempt library/correction, and spending schedule);
changing any evaluation knob implies a new plan id.
An \emph{epoch} (or window) is the declared validity interval of a plan (e.g., a wall-clock period or a beacon-derived id); changing epochs or restart conditions should be explicit.
An \emph{attempt} is one execution of a plan id over a predeclared epoch/window; retries, aborts, and detector changes are attempts and are logged.
A \emph{manifest} is a signed, hash-addressed statement committing code hash, plan id, parameters, and a randomness binding, and pointing to an offline artifact bundle (``available on request'').


\paragraph{Why this matters now.}
Even mature, widely deployed P2P systems continue to see practical deanonymization and routing-layer censorship, underscoring how small, systematic artifacts can collapse anonymity claims if left unmeasured.\cite{heimbach2025_eth_validators,sridhar2024ipfs_censorship}
\paragraph{Related work: anonymous DHTs and lookup leakage.}
Structured-overlay anonymity systems such as Salsa \cite{salsa2006} and AP3 \cite{ap32004} use DHT lookups to locate relays, and follow-on analyses show that lookup mechanisms can leak and be attacked \cite{wangmittalborisov2010,mittalborisov2012}.
Several secure/anonymous DHT designs therefore include \emph{inspection/testing} queries and related defenses \cite{octopus2012,ipersea2014}.
Our contribution is not a new inspection mechanism, but a general evaluation contract: we bound audit--live distinguishability on a published trace schema so that any trace-based metric transfers to deployment up to an explicit, monitored advantage budget.
By design this is orthogonal to runtime drift from a pinned plan: Operational~B asks whether deployed traces still match the declared schedule skeleton, whereas the present paper asks whether the evaluator's probes themselves are detectably special.

\section{Audit awareness as distribution shift}

Audit awareness is a special case of a broader measurement failure mode: if probes are distinguishable from organic traffic,
an adversary (or incentivized intermediary) can treat probes preferentially and bias reported performance.
This problem is well-studied in Internet path-quality measurement \cite{goldberg2008pqm,tariq2009nano} and appears prominently in the censorship evasion arms race,
where censors actively detect and react to measurement/evasion attempts \cite{deresistor2023}.
In anonymous/secure DHT work, the same concern shows up as a design requirement that ``testing'' or ``inspection'' lookups be indistinguishable from ordinary lookups,
otherwise attackers can behave honestly only on probes \cite{octopus2012,ipersea2014,wangmittalborisov2010}.


Let $Z$ be the evaluator's observable trace for a lookup session: timing buckets, hop count, message counts, termination time, etc.
Let $P_{\mathrm{live}}$ be the distribution of $Z$ under deployment traffic and $P_{\mathrm{audit}}$ under the evaluator's probes.
An evaluation-aware system can implement a policy $\pi(z)$ that behaves differently on audit traces, e.g.,
serve probes honestly but bias or leak under live traffic.

\subsection{Advantage and the contract}
Define the audit advantage (TV distance)
\[
\TV(P_{\mathrm{audit}},P_{\mathrm{live}})
\;=\; \sup_{A} \left|P_{\mathrm{audit}}(A)-P_{\mathrm{live}}(A)\right|
\;=\; \frac12 \sup_{\|f\|_{\infty}\le 1} \left|\E_{P_{\mathrm{audit}}}[f(Z)] - \E_{P_{\mathrm{live}}}[f(Z)]\right|.
\]
Equivalently, under equal priors the Bayes-optimal audit detector has accuracy $\mathrm{Acc}^\star$ and advantage $\mathrm{Adv} := 2\,\mathrm{Acc}^\star-1$, and one has $\TV(P_{\mathrm{audit}},P_{\mathrm{live}})=\mathrm{Adv}$.
This is exactly the largest possible audit-live bias for any bounded statistic of the trace.

\paragraph{Remark (hypothesis-testing privacy).}
The advantage contract is a hypothesis-testing constraint on the audit vs.~live trace distributions.
Differential privacy has a closely related ``tradeoff function'' formulation ($f$-DP), of which Gaussian differential privacy (GDP) is a one-parameter family; these languages provide composable accounting for trace randomization when desired.\cite{dong2022gdp}

\begin{theorem}[Advantage contracts bound evaluation bias]\label{thm:adv_contract}
For any measurable $g$ with range in $[0,1]$,
\[
\left|\E_{\mathrm{audit}}[g(Z)] - \E_{\mathrm{live}}[g(Z)]\right| \le \TV(P_{\mathrm{audit}},P_{\mathrm{live}}).
\]
Therefore, if a deployment discipline enforces $\TV(P_{\mathrm{audit}},P_{\mathrm{live}})\le \varepsilon$,
then any bounded evaluation metric computed from traces (risk, leakage proxy, success rate) is biased by at most $\varepsilon$.
\end{theorem}

\paragraph{Remark (range scaling).} If $g$ has range in $[a,b]$ then applying Theorem~\ref{thm:adv_contract} to $(g-a)/(b-a)$ gives $|\E_{\mathrm{audit}}[g]-\E_{\mathrm{live}}[g]|\le (b-a)\,\TV(P_{\mathrm{audit}},P_{\mathrm{live}})$.

\begin{proposition}[Schema-preserving reuse of an advantage contract]\label{prop:reuse_adv}
Let $Z$ be the published trace schema on which an advantage contract establishes $\TV(P_{\mathrm{audit}}^Z,P_{\mathrm{live}}^Z)\le \varepsilon$.
If a later certificate uses a schema $Z' = \phi(Z)$ that is an explicit measurable coarsening of $Z$, then the same contract transfers conservatively:
\[
\TV(P_{\mathrm{audit}}^{Z'},P_{\mathrm{live}}^{Z'}) \le \varepsilon.
\]
If the later certificate changes the schema in any other way, the old contract should not be presented as the same claim; the deployment should publish a new exposure-normal-form id and a fresh advantage contract.
\end{proposition}

\begin{proof}
When $Z'=\phi(Z)$ is a coarsening, the law of $Z'$ is the pushforward of the law of $Z$.
Total variation cannot increase under measurable maps, so the bound follows by data processing.
For non-coarsening schema changes there is, in general, no relation between the old and new distinguishability problems, so reusing the old contract would compare different observables rather than one continuing claim.
\end{proof}

\begin{proposition}[Composing stealth with an audited anonymity certificate]\label{prop:compose}
Let $P_0$ be an ``ideal'' trace distribution for the anonymity objective (as in the earlier series papers), and let $P_{\mathrm{audit}}$ and $P_{\mathrm{live}}$ be the audit and live trace distributions.
If an audited certificate establishes $\TV(P_{\mathrm{audit}},P_0)\le \varepsilon_{\mathrm{cert}}$ and the advantage contract enforces $\TV(P_{\mathrm{audit}},P_{\mathrm{live}})\le \varepsilon_{\mathrm{adv}}$,
then the corresponding live (deployment) trace leakage satisfies
\[
\TV(P_{\mathrm{live}},P_0) \le \varepsilon_{\mathrm{cert}}+\varepsilon_{\mathrm{adv}}.
\]
Thus, an audited trace-level anonymity certificate transfers to deployment up to an additive stealth budget.
\end{proposition}


\begin{table}[t]
\centering
\begin{tabular}{@{}p{0.28\linewidth}p{0.66\linewidth}@{}}
\toprule
Field & Declared public packet content \\
\midrule
Published trace schema id/version & Defines what detectability and bias are measured on (and what detectors are ``covered''). \\
Exposure normal-form id & Stable identifier / schema hash for the exact transcript projection covered by the advantage contract; this is what later certificates must match (or explicitly coarsen) to claim continuity. \\
Epoch / restart rule & Prevents pooled estimates from hiding high-detectability windows; defines the validity interval of $\varepsilon$. \\
Audit selection rule & Makes probes unpredictable before execution and auditable after (e.g., commit + beacon/VRF). \\
Detectability statistic & Specifies the detector/score used for monitoring (e.g., threshold-Youden on a fixed score (a conservative TV lower bound)); enables recomputation. \\
Detector recipe id & Hash commits features/model class/training pipeline to prevent post-hoc detector tuning; changes imply a restart. \\
Lineage relation (optional) & Predecessor contract id plus relation tag (\texttt{same-enf}, \texttt{coarsening}, \texttt{new-interface}) when this contract is compared to an earlier one. \\
Budget $\varepsilon$ & The enforced advantage threshold (TV budget) for transfer of evaluation claims. \\
Publication pointers & Hash/pointer to trace archive plus notarized log entry and inclusion evidence as in the companion release-notary note \cite{series:evalnotary}. \\
\bottomrule
\end{tabular}
\caption{A minimal, artifact-light ``advantage contract certificate'' template.}
\label{tab:adv_manifest}
\end{table}

\subsection{Reader verification checklist (minimal, artifact-light)}
A published advantage contract can remain small while still being referee-checkable.
At minimum, publish (or link to) a compact bundle containing: the trace schema id/version, the exposure normal-form id, the declared epoching/restart rule,
the detector recipe id, the claimed advantage budget $\varepsilon$ (or a confidence bound for it), and the notarized selection/manifests for the audited epochs as in \cite{series:evalnotary}.
That is intentionally the same minimal packet named in the owner stop map above: once those fields are fixed, later notes may certify retries, runtime conformance, or release lineage, but they do not get to redefine the audit-awareness claim exported here.

A referee can then verify, mechanically:
\begin{enumerate}
\item \textbf{Schema coverage.} The published trace schema plausibly covers the detectors the system could use to distinguish audits from live traffic (or else $\varepsilon$ is explicitly a lower bound).
\item \textbf{Lineage discipline.} If the contract is advertised as an update of an earlier one, verify that the exposure normal-form id is unchanged or explicitly declared as a coarsening, as required by Proposition~\ref{prop:reuse_adv}; otherwise treat it as a new contract.
\item \textbf{No post-hoc detector tuning.} The detector recipe id was fixed before evaluation; detectability is reported on held-out epochs/splits under the stated recipe (or changes are logged as restarts).
\item \textbf{Unpredictable probe selection.} Audit selection is commit--reveal/VRF/beacon-derived under the logged commitments, so probes could not be pre-targeted or post-hoc cherry-picked.
\item \textbf{Restart accounting.} Epoch boundaries and restarts are explicit and consistent with the certificate bundle; high-detectability windows are not averaged away.
\end{enumerate}
This is enough for readers to treat the stated $\varepsilon$ as an auditable transfer budget (Theorem~\ref{thm:adv_contract}) without publishing full artifacts inline.

\paragraph{Tiny transfer drill.}
Suppose a deployment publishes the card in Table~\ref{tab:adv_manifest} with exposure normal-form id \texttt{enf:audit-rtt-v2}, an audit/live TV budget $\varepsilon_{\mathrm{adv}}=0.04$, and a trace-level success metric that measured $0.78$ on audited traffic.
Then Theorem~\ref{thm:adv_contract} gives the immediately quotable live interval
\[
0.78 - 0.04 \le \E_{\mathrm{live}}[g] \le 0.78 + 0.04,
\]
so any later terse paper may report the same bounded metric only as lying in $[0.74,0.82]$ unless it widens to a fresh card.
If the audited certificate for the same exposure normal form also established $\TV(P_{\mathrm{audit}},P_0)\le 0.03$, then Proposition~\ref{prop:cert_plus_adv} yields the live leakage bound
\[
\TV(P_{\mathrm{live}},P_0) \le 0.03 + 0.04 = 0.07.
\]
That is enough to quote one audit-awareness packet without reopening evaluator manifests, runtime-plan receipts, or release-lineage packaging.

\paragraph{Interpretation for anonymous DHTs.}
If your anonymity certificate is computed from trace observables, then \emph{stealth} is the missing assumption:
audits must be distributionally close to live traffic, or evaluation-aware nodes can make certificates meaningless.
Moreover, the transfer statement is only as stable as the declared observable: if the covered trace schema changes, the deployment should publish a fresh exposure id and a new advantage contract rather than silently treating the new monitor as ``the same'' one.

Figure~\ref{fig:adv_contract_bias} illustrates the bias bound graphically in a toy model.

\begin{figure}[t]
\centering
\fbox{\begin{minipage}{0.92\linewidth}\centering\small Figure omitted in source-only release. Requested file: \texttt{figures/advantage\_contract\_bias\_bound.png}.\end{minipage}}
\caption{Illustration: bounding trace-level audit advantage controls worst-case bias in bounded metrics.}
\label{fig:adv_contract_bias}
\end{figure}

\section{Monitoring detectability with anytime-valid evidence}
We use anytime-valid evidence (e-values / confidence sequences) to reason about \emph{detectability} of audits under adaptive behavior.
This paper treats those tools as \emph{black-box primitives}: they ensure that a monitor who decides \emph{when} to look (or when to alert) does not accidentally inflate false positives.

\paragraph{Pointer.}
The full ``anytime-valid spending'' story for repeated attempts, optional stopping, and certificate publication is developed in Evaluation Series Paper~1 (Notarized Anonymity Certificates Under Retries).
Here we only need one consequence: if the monitor publishes a nonnegative e-value stream, then adversarially chosen stopping times do not invalidate the declared detection threshold.

\paragraph{What is new here.}
We focus on how an \emph{audit-aware} adversary can shift distributions to evade detection, and on how to pick stealth probes (slice selection) so that detection remains possible without revealing the probe family.

\section{Leakage discovery at scale: slice libraries without log-factor blowup}
Anonymous DHT systems have many potential leakage slices (key-prefix regions, latency regimes, AS paths, churn states, etc.), which creates a many-testing problem.
A practical default is to compute e-values per pre-registered slice and apply the e-BH procedure \cite{wangramdas2022ebh} to control FDR under arbitrary dependence without the strong log-factor penalties of BY-style corrections \cite{benjamini2001by}.
For sequential ``dashboard'' use, there is a subtle filtration issue: repeatedly re-running e-BH under global stopping can fail without additional conditions \cite{stopped_ebh_2025}; see also \cite{xu2024online_evalues}.
For space, the short version omits the full dashboard protocol and toy power comparisons; 

\section{Stealth mechanisms: slice-hiding and verifiable probe selection}
Even if the auditor samples probes from a realistic traffic model, an evaluation-aware overlay can route around \emph{predictable} probes (slice-hiding).
A practical mitigation is to make the set of audited keys/sessions unpredictable until after execution, but verifiable after the fact.

\subsection{Commit--reveal with public randomness}
A minimal pattern is:
\begin{enumerate}
\item commit to an audit pool (e.g., a seed commitment or a hash of candidate probe descriptors),
\item derive the realized audit sample using a public randomness beacon \cite{drand_spec,cloudflare_randomness_beacon_2025,nist_beacon,nist_ir8213_2019} or a VRF \cite{micali1999vrf} (optionally hash-mixing multiple beacons for robustness),
\item reveal and log the selection so anyone can verify it was not post-hoc chosen.
\end{enumerate}
The aim is not secrecy but \emph{unpredictability before execution} and \emph{auditability after execution}.
To keep this verifiable under retries and selective disclosure, treat each initiated selection attempt (including aborts) as a notarized, log-appended manifest entry, and budget restarts with the $\delta$-spending discipline in the companion notarization note \cite{series:evalnotary}.
Similar ``probe resistance'' concerns arise in censorship measurement and in DHT crawling/monitoring infrastructure, where measurement traffic can be fingerprinted and routed around.
IPFS in particular has been repeatedly crawled and monitored in measurement studies, reinforcing that stable audit infrastructure is realistic (and therefore potentially detectable) in deployed Kademlia-family overlays \cite{henningsen2020mapping_ipfs,daniel2022ipfs_churn}.
We treat these as close analogues of audit-aware behavior and cite both the censorship setting and the P2P setting \cite{ensafi2015gfw,deresistor2023,steiner2007kad,wolchok2010dhtcrawl}.
Figure~\ref{fig:slice_hiding} illustrates slice-hiding under predictable audits; the extended version includes a verifiable hidden audit selection (VHAS) toy protocol.

\begin{figure}[t]
\centering
\fbox{\begin{minipage}{0.92\linewidth}\centering\small Figure omitted in source-only release. Requested file: \texttt{figures/slice\_hiding\_audit\_leakage.png}.\end{minipage}}
\caption{Toy slice-hiding effect: predictable audits can be routed around, increasing audit-live bias.}
\label{fig:slice_hiding}
\end{figure}



\section{Padding filters as ``stealth actuators''}
To reduce detectability, the auditor (or system) can apply a padding/actuation filter that maps audit traces toward the live distribution.
This is precisely the type of mechanism treated by the series padding compiler and transport/dual-certificate work \cite{series:schedcompiler,series:knapsacktransport}:
choose a kernel (delay padding, batching, dummy queries) that minimizes overhead subject to a divergence/TV constraint.

In the extended version we include an example of a simple LP-based latency equalization filter from the research archive.



\section{Limitations and open problems}
\paragraph{From AUC to TV.}
AUC is a ranking metric and does not directly upper/lower bound $\TV$ without additional assumptions.
However, any concrete score $S=h(Z)$ yields a conservative \emph{lower bound} on trace-level TV via the best threshold's Youden statistic
(Section 3.2), which can be estimated on held-out data.
Tight, dependence-robust conversions (including confidence sequences for these bounds) remain useful work.
If one is willing to replace classical TV with a smoothed/blurred variant, recent work gives distribution-free upper and lower confidence bounds that may be practical for high-dimensional continuous trace features \cite{hore2026blurredtv}.

\paragraph{Schema coverage.}
The advantage contract only constrains audit--live distinguishability on the \emph{published trace schema} $Z$.
If audit traffic differs on unlogged side channels (e.g., IP ranges, transport fingerprints, lower-layer pacing) or if nodes condition on features omitted from $Z$,
then the contract does not bound evaluation bias for metrics that depend on those omitted features.
For high-assurance claims, the publication should justify why the released schema captures the relevant audit detectors, or treat the advantage budget as a lower bound.

\paragraph{Dependence and churn.}
Anonymous DHT traces are dependent and nonstationary (peer reuse, caching, churn).
Anytime-valid tools remain useful, but audit/live modeling and slice definitions must be revisited under major regime change.
A conservative policy is to treat configuration or membership shifts as restarting a new monitoring epoch.

\paragraph{DHT-native stealth mechanisms.}
The case study is intentionally lightweight; a richer simulator (e.g., Kademlia-like routing with churn and latency models \cite{maymounkov2002kademlia})
would strengthen the engineering story and can be integrated without changing the theoretical contract.

\section{Takeaways}
\begin{itemize}[leftmargin=*]
\item \textbf{Evaluation pipelines are gameable:} if the system can detect audits/probes, it can selectively behave and inflate ``privacy budget'' reports without improving real privacy.
\item \textbf{Advantage contracts separate model from measurement:} define the evaluator's advantage in distinguishing honest from gaming behavior, then design audits to drive that advantage to negligible.
\item \textbf{Stealth and public randomness are first-class:} commit--reveal, randomized reveal order, and notarization patterns reduce cherry-picking and last-revealer attacks in evaluation itself.
\item \textbf{Keep the seam clean:} Eval~1 owns retry-prefix logging, Operational~B owns plan-pinning and epoch conformance of deployed traces, and Release~A owns public lineage; this paper owns only whether audit traffic is itself detectably different from live traffic.
\end{itemize}

\begin{thebibliography}{99}\bibitem{series:schedcompiler}
Anonymous.
\newblock \emph{Scheduling and Compiler Techniques for Metadata-Hiding Overlay Lookups}.
\newblock Series paper (published), 2026.
\newblock Wiki: \texttt{[[2026.01.22 - Mathematics: Scheduling and Compiler Techniques for Metadata-Hiding Overlay Lookups]]}.

\bibitem{series:knapsacktransport}
Anonymous.
\newblock \emph{Optimal Poset-Feasible Padding: Knapsack, Transport, and Dual Certificates for TV Privacy}.
\newblock Series paper (published), 2026.
\newblock Wiki: \texttt{[[2026.01.29 - Mathematics: Optimal Poset-Feasible Padding: Knapsack, Transport, and Dual Certificates for TV Privacy]]}.

\bibitem{series:stoptimeaddendum}
\newblock Stop-Time Padding Addendum (unique-only): Max-Mixture Separations and Tail-Sign Deadline Normal Forms.
\newblock Online manuscript in the author's Anonymity / Anonymous DHT series, 2026.
\newblock Addendum on stopping-time padding pitfalls and normal forms; January 2026

\bibitem{wangramdas2022ebh}
Wang, Ruodu and Ramdas, Aaditya.
\newblock False discovery rate control with e-values.
\newblock Journal of the Royal Statistical Society: Series B, 84(3), 822--852, 2022.

\bibitem{benjamini2001by}
Benjamini, Yoav and Yekutieli, Daniel.
\newblock The control of the false discovery rate in multiple testing under dependency.
\newblock Annals of Statistics, 29(4), 1165--1188, 2001.

\bibitem{maymounkov2002kademlia}
Maymounkov, Petar and Mazieres, David.
\newblock Kademlia: A Peer-to-Peer Information System Based on the XOR Metric.
\newblock International Workshop on Peer-to-Peer Systems (IPTPS), 2002.

\bibitem{micali1999vrf}
Micali, Silvio and Rabin, Michael and Vadhan, Salil.
\newblock Verifiable Random Functions.
\newblock 40th Annual Symposium on Foundations of Computer Science (FOCS), 1999.

\bibitem{drand_spec}
\newblock drand: Distributed Randomness Beacon, Protocol Specification.
\newblock drand documentation, 2026.
\newblock Accessed 2026-02
\newblock \url{https://docs.drand.love/docs/specification/}

\bibitem{nist_beacon}
{NIST}.
\newblock Interoperable Randomness Beacons (NIST CSRC project).
\newblock NIST CSRC project page, 2026.
\newblock Accessed 2026-02
\newblock \url{https://csrc.nist.gov/projects/interoperable-randomness-beacons}

\bibitem{deresistor2023}
Amich, Abderrahmen and Eshete, Birhanu and Yegneswaran, Vinod and Hoang, Nguyen Phong.
\newblock {DeResistor}: Toward {Detection-Resistant} Probing for Evasion of Internet Censorship.
\newblock 32nd USENIX Security Symposium (USENIX Security 23), 2617--2633, 2023.
\newblock \url{https://www.usenix.org/conference/usenixsecurity23/presentation/amich}

\bibitem{ensafi2015gfw}
Ensafi, Roya and Fifield, David and Winter, Philipp and Feamster, Nick and Weaver, Nicholas.
\newblock Examining How the Great Firewall Discovers Hidden Circumvention Servers.
\newblock ACM Internet Measurement Conference (IMC), 2015.
\newblock Active probing and fingerprinting mechanisms; a close analogue of evaluation-aware behavior

\bibitem{steiner2007kad}
Steiner, Moritz and En-Najjary, Taoufik and Biersack, Ernst W..
\newblock A Global View of {KAD}.
\newblock ACM Internet Measurement Conference (IMC), 2007.
\newblock Large-scale measurement of a deployed Kademlia-based DHT; motivates the reality of monitoring/crawling infrastructure
\newblock \url{https://conferences.sigcomm.org/imc/2007/papers/imc60.pdf}

\bibitem{wolchok2010dhtcrawl}
Wolchok, Scott and Halderman, J. Alex.
\newblock Crawling {BitTorrent} {DHT}s for Fun and Profit.
\newblock USENIX Workshop on Offensive Technologies (WOOT), 2010.
\newblock Shows that active crawling creates powerful measurement and adversarial capabilities; relevant to audit detectability and countermeasures
\newblock \url{https://jhalderm.com/pub/papers/dht-woot10.pdf}

\bibitem{daniel2022ipfs_churn}
Daniel, Erik and Tschorsch, Florian.
\newblock Passively Measuring IPFS Churn and Network Size.
\newblock arXiv:2205.14927, 2022.
\newblock Contrasts passive vs active (crawler) visibility in IPFS; accessed 2026-02
\newblock \url{https://arxiv.org/abs/2205.14927}

\bibitem{henningsen2020mapping_ipfs}
Henningsen, Sebastian and Florian, Martin and Rust, Sebastian and Scheuermann, Bj{"o}rn.
\newblock Mapping the InterPlanetary File System.
\newblock IFIP Networking Conference, 2020.
\newblock Large-scale crawl/monitoring study of IPFS' Kademlia DHT; accessed 2026-02
\newblock \url{https://dl.ifip.org/db/conf/networking/networking2020/1570620406.pdf}

\bibitem{xu2024online_evalues}
Xu, Ziyu and Liang, Tengyuan and Ramdas, Aaditya.
\newblock Online multiple testing with e-values.
\newblock Proceedings of The 27th International Conference on Artificial Intelligence and Statistics (AISTATS), 2024.
\newblock Doubly-sequential frameworks combining sequential testing with online multiple testing
\newblock \url{https://proceedings.mlr.press/v238/xu24a/xu24a.pdf}

\bibitem{series:evalnotary}
\newblock Notarized Anonymity Certificates Under Retries: Grinding-Resistant Evaluation via Transparency Logs, Public Randomness, and $\delta$-Spending.
\newblock Online manuscript in the author's Anonymity / Anonymous DHT series, 2026.

\bibitem{dong2022gdp}
Dong, Jinshuo and Roth, Aaron and Su, Weijie J..
\newblock Gaussian Differential Privacy.
\newblock Journal of the Royal Statistical Society: Series B (Statistical Methodology), 84(1), 3--37, 2022.
\newblock DOI: 10.1111/rssb.12455.
\newblock Introduces $f$-DP (hypothesis-testing privacy) and GDP; provides interpretable, composable accounting for distributional indistinguishability

\bibitem{stopped_ebh_2025}
Wang, Hongjian and Dandapanthula, Sanjit and Ramdas, Aaditya.
\newblock Anytime-valid FDR control with the stopped e-BH procedure.
\newblock arXiv:2502.08539, 2025.
\newblock Accessed 2026-02
\newblock \url{https://arxiv.org/abs/2502.08539}

\bibitem{cloudflare_randomness_beacon_2025}
\newblock Cloudflare Randomness Beacon (drand) documentation.
\newblock Cloudflare developer documentation, 2025.
\newblock Accessed 2026-02
\newblock \url{https://developers.cloudflare.com/randomness-beacon/}

\bibitem{nist_ir8213_2019}
Kelsey, John.
\newblock {NIST IR 8213}: A Reference for Randomness Beacons (Format and Protocol Version 2).
\newblock NIST CSRC, 2019.
\newblock Initial public draft published May 2019; accessed 2026-02
\newblock \url{https://csrc.nist.gov/pubs/ir/8213/ipd}

\bibitem{hore2026blurredtv}
Hore, Rohan and Foygel Barber, Rina.
\newblock Distribution-free two-sample testing with blurred total variation distance.
\newblock 2026.
\newblock Distribution-free upper/lower confidence bounds for a smoothed (blurred) relaxation of TV, enabling nonparametric inference for continuous/high-dimensional traces
\newblock \url{https://arxiv.org/abs/2602.05862}

\bibitem{series:traceifaces}
Anonymous.
\newblock Trace interfaces for anonymous DHTs.
\newblock Synthesis~3, 2026. (This archive.)

\bibitem{dwork2015reusableholdout}
Dwork, Cynthia and Feldman, Vitaly and Hardt, Moritz and Pitassi, Toniann and Reingold, Omer and Roth, Aaron.
\newblock The reusable holdout: Preserving validity in adaptive data analysis.
\newblock Science, 349(6248), 636--638, 2015.
\newblock DOI: 10.1126/science.aaa9375.
\newblock \url{https://www.science.org/doi/10.1126/science.aaa9375}

\bibitem{dwork2014adaptive}
Dwork, Cynthia and Feldman, Vitaly and Hardt, Moritz and Pitassi, Toniann and Reingold, Omer and Roth, Aaron.
\newblock Preserving Statistical Validity in Adaptive Data Analysis.
\newblock 2014.
\newblock \url{https://arxiv.org/abs/1411.2664}

\bibitem{octopus2012}
Wang, Qiyan and Borisov, Nikita.
\newblock Octopus: A Secure and Anonymous DHT Lookup.
\newblock 2012 IEEE 32nd International Conference on Distributed Computing Systems (ICDCS), 325--334, 2012.
\newblock DOI: 10.1109/ICDCS.2012.78.
\newblock Author-hosted PDF; accessed 2026-02
\newblock \url{https://nikita.ca/papers/octopus-icdcs12.pdf}

\bibitem{heimbach2025_eth_validators}
Heimbach, Lioba and Vonlanthen, Yann and Villacis, Juan and Kiffer, Lucianna and Wattenhofer, Roger.
\newblock Deanonymizing Ethereum Validators: The P2P Network Has a Privacy Issue.
\newblock 34th USENIX Security Symposium (USENIX Security 25), 2025.
\newblock Large-scale deanonymization in a deployed P2P network; accessed 2026-02
\newblock \url{https://www.usenix.org/system/files/usenixsecurity25-heimbach.pdf}

\bibitem{ipersea2014}
Al-Ameen, Mahdi Nasrullah and Wright, Matthew.
\newblock iPersea: The Improved Persea with Sybil Detection Mechanism.
\newblock arXiv:1412.6883, 2014.
\newblock Introduces ``inspection lookup'' for Sybil detection in a social DHT
\newblock \url{https://arxiv.org/abs/1412.6883}

\bibitem{sridhar2024ipfs_censorship}
Sridhar, Srivatsan and Ascigil, Onur and Keizer, Navin and Genon, Fran{\c{c}}ois and Pierre, S{\'e}bastien and Psaras, Yiannis and Riviere, Etienne and Kr{\'o}l, Micha{\l}.
\newblock Content Censorship in the InterPlanetary File System.
\newblock Network and Distributed System Security Symposium (NDSS), 2024.
\newblock DOI: 10.14722/ndss.2024.23153.
\newblock Demonstrates low-cost DHT-layer censorship in IPFS and mitigation; accessed 2026-02
\newblock \url{https://www.ndss-symposium.org/wp-content/uploads/2024-153-paper.pdf}

\bibitem{salsa2006}
Nambiar, Arjun and Wright, Matthew.
\newblock Salsa: A Structured Approach to Large-Scale Anonymity.
\newblock Proceedings of the 13th ACM Conference on Computer and Communications Security (CCS), 2006.
\newblock DOI: 10.1145/1180405.1180409.
\newblock Structured-overlay anonymity that uses DHT lookups to select relays
\newblock \url{https://www.freehaven.net/anonbib/cache/Salsa.pdf}

\bibitem{ap32004}
Mislove, Alan and Oberoi, Gaurav and Post, Ansley and Reis, Charles and Druschel, Peter and Wallach, Dan S..
\newblock AP3: Cooperative, Decentralized Anonymous Communication.
\newblock Proceedings of the 11th ACM SIGOPS European Workshop: Beyond the PC, 2004.
\newblock DOI: 10.1145/1133572.1133578.
\newblock Anonymous overlay with cooperative relay selection; uses DHT-like lookups
\newblock \url{https://www.cs.rice.edu/~dwallach/pub/ap3-sigops04.html}

\bibitem{wangmittalborisov2010}
Wang, Qiyan and Mittal, Prateek and Borisov, Nikita.
\newblock In Search of an Anonymous and Secure Lookup: Attacks on Structured Peer-to-Peer Anonymous Communication Systems.
\newblock Proceedings of the 17th ACM Conference on Computer and Communications Security (CCS), 308--318, 2010.
\newblock DOI: 10.1145/1866307.1866343.
\newblock \url{https://www.freehaven.net/anonbib/cache/ccs10-lookup.pdf}

\bibitem{mittalborisov2012}
Mittal, Prateek and Borisov, Nikita.
\newblock Information Leaks in Structured Peer-to-Peer Anonymous Communication Systems.
\newblock ACM Transactions on Information and System Security, 15(1), 5:1--5:28, 2012.
\newblock DOI: 10.1145/2133375.2133380.
\newblock \url{https://www.princeton.edu/~pmittal/publications/information-leaks-tissec.pdf}

\bibitem{goldberg2008pqm}
Goldberg, Sharon and Rexford, Jennifer and Jaggard, Aaron D. and Ramachandran, Anitha and Turner, J. S. and Vyas, Prateek.
\newblock Path-Quality Monitoring in the Presence of Adversaries.
\newblock ACM SIGMETRICS, 2008.
\newblock Shows that active probing can be gamed when probes are treated preferentially
\newblock \url{https://www.cs.princeton.edu/~jrex/papers/sigmetrics08-pqm.pdf}

\bibitem{tariq2009nano}
Tariq, Mukarram Bin and Motiwala, Murtaza and Feamster, Nick and Ammar, Mostafa.
\newblock Detecting Network Neutrality Violations with Causal Inference.
\newblock ACM CoNEXT, 2009.
\newblock Notes that active probes are detectable and can be preferentially treated; accessed 2026-02
\newblock \url{https://www.gtnoise.net/papers/2009/tariq_nano_conext09.pdf}

\bibitem{series:workedexample}
Anonymous.
\newblock Worked Example: Instantiating a Tiered Reader-Privacy Receipt.
\newblock Synthesis~17, The-One-True-Archive (2026).

\bibitem{series:pairedterminalcutoverrules}
Anonymous.
\newblock Paired Terminal Cutover Rules for Stable Review Spines: When Later Terse Papers Must Stop, Widen, or Reopen.
\newblock Series synthesis note (Synthesis~75), 2026. (This archive.)

\bibitem{series:artifactpolicy}
Anonymous.
\newblock Artifact policy (source-only archive; artifacts available on request).
\newblock Series note, 2026.

\bibitem{series:relatedwork}
Anonymous.
\newblock Related Work Map and Pointer Index for the One-True-Archive.
\newblock Series synthesis note (Synthesis~7), 2026. (This archive.)

\end{thebibliography}

\end{document}